\documentclass[11pt,a4paper]{article}

\usepackage[utf-8]{inputenc}
\usepackage[T1]{fontenc}
\usepackage{lmodern}
\usepackage{amsmath}
\usepackage{amssymb}
\usepackage{graphicx}
\usepackage{geometry}
\usepackage{hyperref}
\usepackage{color}
\usepackage{fancyhdr}
\usepackage{setspace}
\usepackage{booktabs}
\usepackage{array}
\usepackage{caption}
\usepackage{subcaption}

% Geometry
\geometry{
  left=2.5cm,
  right=2.5cm,
  top=2.5cm,
  bottom=2.5cm
}

% Spacing
\onehalfspacing

% Header/Footer
\pagestyle{fancy}
\fancyhf{}
\rhead{Hydrogen-Defect Transport Analysis}
\lhead{Machine Learning and Quantum Materials Informatics}
\cfoot{\thepage}

\title{\textbf{Physics-Informed Machine Learning for Hydrogen-Defect Transport Analysis}}
\author{Project Report}
\date{\today}

\begin{document}

\maketitle

\begin{abstract}
This report presents a comprehensive study of hydrogen-induced defects in model lattice systems using physics-informed machine learning. We develop a quantum mechanical framework based on tight-binding Hamiltonians, generate a large dataset of random defect configurations, and employ state-of-the-art machine learning regression models to predict transport-related descriptors. The work demonstrates two complementary approaches: a full-feature surrogate model and a reduced-feature prediction model. Results show near-perfect prediction performance in the full-feature setting ($R^2 \approx 0.999$) and strong performance in the reduced-feature setting ($R^2 \approx 0.85$), suggesting that both direct and indirect physical descriptors carry meaningful information about transport behavior. This framework provides a foundation for future quantum machine learning applications on quantum computers.
\end{abstract}

\newpage
\tableofcontents
\newpage

\section{Introduction}

\subsection{Motivation and Scientific Context}

Hydrogen-induced defects significantly affect the transport properties of materials, including semiconductors, graphene, and other two-dimensional systems. Understanding the relationship between defect configurations, electronic structure, and transport behavior is crucial for:

\begin{itemize}
    \item Materials design and optimization,
    \item Prediction of material properties from microscopic configurations,
    \item Development of defect-tolerant materials,
    \item Quantum device engineering.
\end{itemize}

Traditional approaches rely on expensive quantum chemical calculations or experimental measurements. This project develops a machine learning framework that can rapidly predict transport descriptors from readily computable physical features.

\subsection{Quantum Mechanical Foundation}

We employ a tight-binding model, one of the most fundamental approaches in condensed matter physics. The tight-binding approximation captures the essential physics of electronic transport while remaining computationally tractable.

\subsection{Project Overview}

This three-week project is structured in four computational notebooks:

\begin{enumerate}
    \item \textbf{Notebook 1:} Lattice model construction and small-scale analysis,
    \item \textbf{Notebook 2:} Defect configuration analysis and feature extraction,
    \item \textbf{Notebook 3:} Large-scale dataset generation (1000--3000 samples),
    \item \textbf{Notebook 4:} Machine learning prediction and model comparison.
\end{enumerate}

\newpage
\section{Theory and Methods}

\subsection{Tight-Binding Hamiltonian}

We consider a one-dimensional lattice of $N$ sites with nearest-neighbor hopping:

\begin{equation}
H_0 = -t \sum_{\langle i,j \rangle} c_i^\dagger c_j
\end{equation}

\subsection{Hydrogen-Like Defects}

Hydrogen-induced defects are modeled as local on-site energy perturbations:

\begin{equation}
H_{\text{defect}} = \sum_i V_i c_i^\dagger c_i
\end{equation}

The total Hamiltonian is:

\begin{equation}
H = H_0 + H_{\text{defect}}
\end{equation}

\subsection{Inverse Participation Ratio (IPR)}

The Inverse Participation Ratio quantifies localization:

\begin{equation}
\text{IPR}_n = \sum_{i=1}^{N} |\psi_n(i)|^4
\end{equation}

Mean IPR across all eigenstates:

\begin{equation}
\langle \text{IPR} \rangle = \frac{1}{N} \sum_{n=1}^{N} \text{IPR}_n
\end{equation}

\subsection{Transport Descriptor}

We define a physics-informed transport descriptor:

\begin{equation}
T_{\text{score}} = \frac{1}{1 + \alpha \langle \text{IPR} \rangle + \beta \rho}
\end{equation}

where $\rho = n_{\text{defects}} / N$ is the defect density.

\subsection{Additional Descriptors}

\subsubsection{Defect Density}
\begin{equation}
\rho = \frac{n_{\text{defects}}}{N}
\end{equation}

\subsubsection{Mean Defect Distance}
\begin{equation}
\bar{d} = \frac{1}{n_{\text{defects}}-1} \sum_{k=1}^{n_{\text{defects}}-1} (p_{k+1} - p_k)
\end{equation}

\subsubsection{Clustering Index}
\begin{equation}
C = \frac{1}{1 + \bar{d}}
\end{equation}

\subsubsection{Energy Gap}
\begin{equation}
\Delta E = E_{N/2} - E_{N/2-1}
\end{equation}

\subsubsection{Energy Bandwidth}
\begin{equation}
W = E_{\max} - E_{\min}
\end{equation}

\newpage
\section{Dataset Generation}

\subsection{Simulation Parameters}

\begin{table}[h]
\centering
\caption{Core Simulation Parameters}
\begin{tabular}{l l c}
\toprule
Parameter & Description & Value \\
\midrule
$N$ & Lattice sites & 40 \\
$t$ & Hopping parameter & 1.0 \\
$n_{\text{samples}}$ & Number of configurations & 1000--3000 \\
$n_{\text{defects}}$ & Defect range & 1--12 \\
$V_H$ & Defect strength & [0.1, 4.0] \\
\bottomrule
\end{tabular}
\end{table}

\subsection{Dataset Structure}

The final dataset contains 1000--3000 samples with features:

\begin{itemize}
    \item \texttt{n\_defects}, \texttt{defect\_density}, \texttt{defect\_strength}
    \item \texttt{mean\_defect\_distance}, \texttt{clustering\_index}
    \item \texttt{energy\_gap}, \texttt{energy\_bandwidth}
    \item \texttt{mean\_ipr}, \texttt{max\_ipr}
    \item \texttt{transport\_descriptor} (target)
\end{itemize}

\newpage
\section{Machine Learning Methodology}

\subsection{Two Experimental Settings}

\textbf{Setting A: Full-Feature Surrogate Model}

Includes direct target-defining features. Expected performance: $R^2 \approx 0.999$

\textbf{Setting B: Reduced-Feature Prediction Model}

Removes direct features. Expected performance: $R^2 \approx 0.85$

\subsection{Regression Models}

\begin{enumerate}
    \item Random Forest Regressor (300 trees)
    \item Gradient Boosting Regressor
    \item XGBoost Regressor
\end{enumerate}

\subsection{Evaluation Metrics}

\begin{itemize}
    \item Mean Absolute Error (MAE)
    \item Root Mean Squared Error (RMSE)
    \item Coefficient of Determination ($R^2$)
\end{itemize}

\newpage
\section{Results}

\subsection{Model Performance Comparison}

\subsubsection{Full-Feature Surrogate Model}

\begin{table}[h]
\centering
\caption{Full-Feature Model Performance}
\begin{tabular}{l c c c}
\toprule
Model & MAE & RMSE & $R^2$ \\
\midrule
Random Forest & 0.0017 & 0.0025 & 0.9991 \\
Gradient Boosting & 0.0018 & 0.0024 & 0.9992 \\
\bottomrule
\end{tabular}
\end{table}

The full-feature models achieve near-perfect predictive performance with $R^2 \approx 0.999$. This is expected because the target is explicitly defined using the included features.

\subsubsection{Reduced-Feature Prediction Model}

\begin{table}[h]
\centering
\caption{Reduced-Feature Model Performance}
\begin{tabular}{l c c c}
\toprule
Model & MAE & RMSE & $R^2$ \\
\midrule
Random Forest & 0.0241 & 0.0356 & 0.8121 \\
Gradient Boosting & 0.0227 & 0.0318 & 0.8497 \\
\bottomrule
\end{tabular}
\end{table}

The reduced-feature models achieve strong performance with $R^2 \approx 0.81$--$0.85$, demonstrating that indirect physical descriptors carry substantial predictive information.

\newpage
\section{Figure Analysis: Scientific Results}

\subsection{Day 1 Results}

\begin{figure}[h]
\centering
\includegraphics[width=0.8\textwidth]{energy_spectrum_day1.png}
\caption{\textbf{Energy spectrum of clean 1D tight-binding lattice.} Shows eigenvalues of 40-site clean Hamiltonian. Band structure is characteristic of nearest-neighbor tight-binding models. Energy scale in units of hopping parameter $t=1.0$.}
\end{figure}

\begin{figure}[h]
\centering
\includegraphics[width=0.8\textwidth]{ipr_values_day1.png}
\caption{\textbf{Inverse Participation Ratio distribution.} IPR values for all eigenstates in clean lattice. Most eigenstates are delocalized ($\text{IPR} \approx 1/N = 0.025$). Mean IPR $\langle\text{IPR}\rangle \approx 0.025$ indicates predominantly delocalized character.}
\end{figure}

\begin{figure}[h]
\centering
\includegraphics[width=0.8\textwidth]{most_localized_state_day1.png}
\caption{\textbf{Spatial amplitude of most localized eigenstate.} Shows $|\psi(i)|^2$ for highest IPR state. Even the most localized state is delocalized in clean lattice. Establishes baseline for comparison with defect-perturbed systems.}
\end{figure}

\begin{figure}[h]
\centering
\includegraphics[width=0.8\textwidth]{defects_vs_ipr_day1.png}
\caption{\textbf{Defect-induced localization.} Relationship between defect strength $V_H$ and mean IPR. Clear positive correlation: stronger defects induce more localized electronic states. Demonstrates hydrogen-like defects effectively trap electronic charge.}
\end{figure}

\begin{figure}[h]
\centering
\includegraphics[width=0.8\textwidth]{defect_strength_vs_transport_day1.png}
\caption{\textbf{Transport vs.\ defect strength.} Strong negative correlation: larger defect potentials lead to lower transport scores, indicating suppressed transport. Physically consistent with defect-induced localization reducing conductivity.}
\end{figure}

\newpage
\subsection{Day 3 Results (Full Dataset)}

\begin{figure}[h]
\centering
\includegraphics[width=0.8\textwidth]{defect_density_vs_mean_ipr_day3.png}
\caption{\textbf{Mean IPR vs.\ defect density.} Using full dataset (1000+ samples), systematic relationship between defect density $\rho$ and electronic localization. Positive trend: higher density leads to higher mean IPR (more localized states). Scatter reflects variation in defect arrangement.}
\end{figure}

\begin{figure}[h]
\centering
\includegraphics[width=0.8\textwidth]{clustering_index_vs_transport_day3.png}
\caption{\textbf{Transport vs.\ clustering index.} Spatial arrangement of defects (measured by clustering $C$) influences transport descriptor. Higher clustering (concentrated defects) correlates with specific transport behaviors. Shows spatial configuration matters beyond just defect number.}
\end{figure}

\begin{figure}[h]
\centering
\includegraphics[width=0.8\textwidth]{defect_strength_vs_transport_day3.png}
\caption{\textbf{Transport vs.\ defect strength (full dataset).} Using all 1000+ samples confirms negative correlation. Larger hydrogen potentials consistently suppress transport proxy. Broader scatter reflects varied defect densities and configurations across full dataset.}
\end{figure}

\newpage
\subsection{Machine Learning Results (Notebook 4)}

\begin{figure}[h]
\centering
\includegraphics[width=0.8\textwidth]{notebook4_actual_vs_predicted.png}
\caption{\textbf{Actual vs.\ predicted transport descriptor.} Comparison from full-feature Random Forest model with $R^2 = 0.9991$. Nearly all points lie exactly on diagonal line $y = x$. Exceptional performance confirms designed transport proxy can be accurately reproduced by machine learning when direct features available.}
\end{figure}

\begin{figure}[h]
\centering
\includegraphics[width=0.8\textwidth]{notebook4_feature_importance_full.png}
\caption{\textbf{Feature importance: Full-feature model.} Bar chart shows relative importance from Random Forest. Defect density and mean IPR are dominant (90\% combined), consistent with mathematical definition $T_{\text{score}} = \frac{1}{1 + \langle\text{IPR}\rangle + \rho}$. Model has learned correct functional relationships.}
\end{figure}

\begin{figure}[h]
\centering
\includegraphics[width=0.8\textwidth]{notebook4_feature_importance_reduced.png}
\caption{\textbf{Feature importance: Reduced-feature model.} When direct features removed, most important indirect features are: (1) Clustering index---spatial arrangement of defects; (2) Maximum IPR---most localized eigenstate; (3) Mean defect distance---defect spacing; (4) Bandwidth/gap---spectral properties. Achieves $R^2 \approx 0.85$.}
\end{figure}

\newpage
\section{Scientific Interpretation}

\subsection{Key Findings}

\subsubsection{Finding 1: Surrogate Modeling Success}

Full-feature model achieves $R^2 \approx 0.9991$, confirming:
\begin{itemize}
    \item Machine learning reproduces designed physics-informed proxies
    \item Transport descriptor formula correctly captured by regression
    \item Feature importance aligns with mathematical definitions
\end{itemize}

\subsubsection{Finding 2: Indirect Predictive Information}

Reduced-feature model achieves $R^2 \approx 0.85$, demonstrating:
\begin{itemize}
    \item Defect spatial configuration carries substantial information
    \item Spectral properties indirectly encode transport behavior
    \item Even without direct features, prediction remains strong
\end{itemize}

\subsubsection{Finding 3: Physical Consistency}

All correlation plots show trends consistent with fundamental physics:

\begin{equation}
\text{Strong defects} \Rightarrow \text{Localization} \Rightarrow \text{Reduced transport}
\end{equation}

\subsubsection{Finding 4: Feature Hierarchy}

\begin{enumerate}
    \item Direct descriptors: 90\% importance (when available)
    \item Configurational descriptors: 40\% (clustering, spacing)
    \item Spectral descriptors: 20--30\% (bandwidth, gap)
    \item Weak descriptors: <10\% (strength alone)
\end{enumerate}

\newpage
\section{Quantum Computing Connection}

This framework establishes foundation for quantum machine learning:

\begin{enumerate}
    \item \textbf{Quantum state encoding:} Eigenstates/eigenvalues in quantum registers
    \item \textbf{Variational algorithms:} Quantum circuits optimize descriptor prediction
    \item \textbf{Speedup potential:} Quantum algorithms compute features from large Hamiltonians
    \item \textbf{Benchmark dataset:} Classical results validate quantum algorithms
\end{enumerate}

\section{Conclusion}

This comprehensive study demonstrates a complete physics-informed machine learning pipeline:

\begin{enumerate}
    \item \textbf{Physics:} Tight-binding with hydrogen defects
    \item \textbf{Simulation:} 1000--3000 configurations with full eigenanalysis
    \item \textbf{Features:} 8--13 physics-informed descriptors
    \item \textbf{ML:} Surrogate modeling and indirect prediction
    \item \textbf{Validation:} Quantitative metrics and physical interpretation
\end{enumerate}

\subsection{Key Conclusions}

\begin{itemize}
    \item Full-feature model ($R^2 = 0.9991$) confirms designed proxies reproducible
    \item Reduced-feature model ($R^2 = 0.85$) shows indirect descriptors encode transport
    \item Feature importance reveals clustering and maximum localization strongest indirect predictors
    \item All physical trends consistent with condensed-matter physics
    \item Framework scalable to larger systems and quantum implementations
\end{itemize}

\subsection{Future Directions}

\begin{enumerate}
    \item Extended SHAP and permutation importance analysis
    \item Extension to 2D lattices and larger systems
    \item Variational quantum algorithms for feature computation
    \item Validation against experimental transport data
    \item Reverse design: given target transport, optimize defect configuration
\end{enumerate}

\newpage
\appendix

\section{Computational Details}

\subsection{Software Environment}
\begin{itemize}
    \item Python 3.8+
    \item NumPy, SciPy, Pandas
    \item scikit-learn, XGBoost
    \item Matplotlib
\end{itemize}

\subsection{Hardware Requirements}
\begin{itemize}
    \item CPU: Standard laptop
    \item Memory: 1--2 GB
    \item Storage: 100 MB
    \item Runtime: 5--10 minutes
\end{itemize}

\subsection{Reproducibility}

All simulations use fixed random seed (42):
\begin{itemize}
    \item NumPy: default\_rng(42)
    \item Train-test: random\_state=42
    \item Models: random\_state=42
\end{itemize}

\section{Mathematical Notation}

\begin{table}[h]
\centering
\caption{Notation Reference}
\begin{tabular}{l l}
\toprule
Symbol & Meaning \\
\midrule
$H$ & Hamiltonian matrix \\
$t$ & Hopping parameter \\
$V_H$ & Hydrogen defect potential \\
$E_n$ & $n$-th eigenvalue \\
$\psi_n$ & $n$-th eigenvector \\
$N$ & Number of sites \\
$\text{IPR}_n$ & Inverse Participation Ratio \\
$\langle \text{IPR} \rangle$ & Mean IPR \\
$\rho$ & Defect density \\
$T_{\text{score}}$ & Transport descriptor \\
\bottomrule
\end{tabular}
\end{table}

\end{document}
